\section{Continuous formulas with integration binders}\label{sec:logic}

The earlier compact-model packets already give quantitative meanings to
infimum and supremum formulas \cite{repo-abstract,repo-continuous-two}.
We extend their interface with integration. Continuous integration logic
is an established subject \cite{bagheri-pourmahdian}; the result here is a
finite-evaluation theorem with explicit hypotheses, rather than a proof of
its unrestricted compactness or ultraproduct theorems.

\subsection{Compact probability sorts and a constructive quadrature}

A \emph{compact probability sort} consists of a presented metric space
with a rational diameter bound, PR finite $h$-nets for rational $h>0$,
and a presented probability integral. Each named distance-to-center
function is supplied with compatible $L^1$ observations. The integral is
positive, takes $1$ to $1$, and has the contraction property of the measure
chapter. Compatibility is substantive: arbitrary compact and measure
presentations on unrelated carriers do not constitute one sort.

Here is an explicit sufficient construction of quadrature from this data.
It avoids having to decide the boundaries of nearest-center cells.

\begin{lemma}[Tent quadrature with rational masses]\label{lem:tent-quadrature}
Let $q_1,\ldots,q_m$ be an $h/2$-net, $m\ge1$, and set
\[
 t_i(x)=(h-d(x,q_i))_+,\quad S(x)=\sum_it_i(x),\quad
 \phi_i(x)=\frac{t_i(x)}{S(x)}.
\]
Then $S\ge h/2$, $\phi_i\ge0$, $\sum_i\phi_i=1$, and
$\phi_i(x)>0$ implies $d(x,q_i)<h$.
Let $m_i=I(\phi_i)$. If rational $a_i$ satisfy $|a_i-m_i|\le\eta$
and $m\eta\le1/2$, define
\[
 u_i=\max(0,a_i),\qquad U=\sum_iu_i,\qquad w_i=u_i/U.
\]
Then $U\ge1/2$, the $w_i$ form a rational probability vector, and
$\sum_i|w_i-m_i|\le2m\eta$.
For every $L$-Lipschitz $g$ with $|g|\le B$,
\begin{equation}\label{eq:quadrature}
 \left|I(g)-\sum_iw_i g(q_i)\right|
 \le Lh+2Bm\eta.
\end{equation}
All numerical choices are PR from the net, the integral names, and the
requested error.
\end{lemma}
\begin{proof}
Some center lies at distance at most $h/2$ from every $x$, so one tent is
at least $h/2$. This proves the positive denominator bound and the
partition identities. The reciprocal is evaluated using the globally
defined guarded function $r(s)=1/\max(h/2,s)$. It is bounded by $2/h$
and is $4/h^2$-Lipschitz. Approximate tents are clipped to $[0,h]$
before multiplication. If their $L^1$ errors are $e_j$, the resulting
error in $\phi_i$ is at most $(2/h)e_i+(4/h)\sum_je_j$. Thus intermediate
approximations need not themselves satisfy the exact lower bound on $S$.
The bounded product rule gives named $L^1$ functions $\phi_i$.

Positivity and unit mass imply $m_i\ge0$ and $\sum_im_i=1$.
Clipping $a_i$ at zero does not increase its distance from $m_i$, hence
$D=\sum_i|u_i-m_i|\le m\eta$. Thus $|U-1|\le D$ and $U\ge1/2$.
Moreover
\[
 \sum_i|u_i/U-m_i|
 \le\sum_i|u_i/U-u_i|+D
 =|1-U|+D\le2D.
\]
If $g$ is initially only a bounded Lipschitz point-evaluator, the tent
interpolants $g_h=\sum_i g(q_i)\phi_i$ first construct its $L^1$ name.
Pointwise $|g-g_h|\le Lh$; therefore any two such interpolants differ
in $L^1$ by at most the sum of their $Lh$ bounds, and a dyadic schedule
gives a regular name. This proves the needed compatibility between point
evaluation and integration rather than assuming it for every new atom.
For exact masses, use the partition identity inside the integral:
\[
 \left|I(g)-\sum_im_i g(q_i)\right|
 \le I\left(\sum_i\phi_i(x)|g(x)-g(q_i)|\right)\le Lh.
\]
The last inequality uses the support condition of each tent. Perturbing
the masses costs at most $B\sum_i|w_i-m_i|$, proving
\eqref{eq:quadrature}. Evaluating the finitely many mass names with
$\eta\le\min(1/(2m),\eps/(4Bm))$ and taking $h\le\eps/(2L)$ gives
a total error at most $\eps$, with the evident separate cases $B=0$ or
$L=0$. These are positive rational divisions and bounded loops.
\end{proof}

The tent construction itself is already present in the repository's
function-space work. The new use is to convert its support and denominator
certificates into probability quadratures with normalized rational masses.
No exact decision of a named real inequality is inserted into the algorithm.

\subsection{Syntax, moduli, and recursive evaluation}

Fix a finite many-sorted signature. Each atomic predicate is a named
bounded real-valued function of its arguments, with rational coordinate
Lipschitz bounds $L_{R,j}$ and a jointly PR evaluator on dense arguments.
Function symbols can be admitted by first computing their term moduli and
expanding them into the atomic evaluator; their evaluation error is charged
there. A formula is a finite expression generated by:
\begin{enumerate}
\item rational constants and atomic predicates;
\item rational, bounded Lipschitz connectives, including complement,
      minimum, maximum, affine averaging, and products on $[0,1]$;
\item $\inf_x\varphi$, $\sup_x\varphi$ over a specified compact sort;
\item $\mathsf E_x\varphi$, integrating over that sort's fixed probability
      measure, independent of the remaining free parameters.
\end{enumerate}
The bound-variable positions and sorts are part of the syntax. We do not
permit an unannounced change of the measure with the environment. A
parameter-dependent kernel binder requires the additional kernel-variation
contract from Section~\ref{sec:operators}.

Range bounds $B_\varphi$ and Lipschitz constants for each free variable are
computed structurally. For a connective with
$|C(a)-C(b)|\le\sum_jc_j|a_j-b_j|$, use the corresponding nonnegative
sum of the input constants. For minimum and maximum one can use the
maximum of coordinate constants. A binder drops the bound coordinate
and leaves the other constants unchanged. Integration has mass one;
infimum and supremum preserve a uniform parameter bound. The connective
range certificate supplies the next $B_\varphi$.

For a requested precision the finite evaluator replaces bound sorts by
chosen finite nets, infimum and supremum by finite minimum and maximum,
and expectation by the chosen rational weighted sum. Before using a
connective whose Lipschitz law is certified only on specified range
intervals, clip each approximate argument to its certified interval.
Clip each output to its proved range as well. Such clipping cannot
increase the error to the true value; it ensures, for example, that
the product rule on $[0,1]$ is never applied to an out-of-range sample.
The evaluator retains the syntactic nesting order. All tables are finite and all loops have their
explicit lengths. A fixed formula has a PR evaluator; a numerical syntax
code can be interpreted by bounded syntax recursion with its program
clocks, but this does not provide a universal evaluator for arbitrary PR
symbols.

\begin{theorem}[Mixed quantifier evaluation]\label{thm:mixed-evaluation}
With the preceding data, every fixed well-sorted formula defines a named
continuous function of its free variables. Suppose its atomic tables have
uniform error budgets, and each integration node uses a quadrature with
radius $h$ and mass error at most $\tau$. A sound numerical error recursion
is as follows:
\begin{align*}
 e_R&=\text{the supplied atomic evaluation error},\\
 e_{C(\varphi_1,\ldots,\varphi_s)}&=\sum_jc_je_{\varphi_j},\\
 e_{\min(\varphi,\psi)},e_{\max(\varphi,\psi)}
     &\le\max(e_\varphi,e_\psi),\\
 e_{\inf_x\varphi},e_{\sup_x\varphi}
     &\le e_\varphi+L_{\varphi,x}h,\\
 e_{\mathsf E_x\varphi}
     &\le e_\varphi+L_{\varphi,x}h+B_\varphi\tau.
\end{align*}
Here $h$ in an extremum rule is a covering radius. The output error is at
most the budget computed at the root. A PR precision schedule makes that
budget at most $2^{-p}$.
\end{theorem}
\begin{proof}
Induct on the fixed syntax. Constants and atoms are their supplied names.
The connective step is exactly its Lipschitz inequality, using the
inductively obtained range bounds. It also proves the required free-variable
moduli, so the integrands in the next steps remain uniformly continuous.

For a supremum of a function $g$ with modulus $L$, let
$M=\max_i g(q_i)$. Every $x$ is within $h$ of some $q_i$, and hence
$g(x)\le M+Lh$; conversely every $g(q_i)$ is a value of $g$. Thus the
finite maximum approximates the supremum within $Lh$. Perturbing all
table values by at most $e$ changes the maximum by at most $e$.
The infimum proof reverses the inequalities. This argument can also
\emph{construct} the supremum name without assuming its existence:
two finite nets with errors $(h,e)$ and $(h',e')$ give maxima differing by
at most $L(h+h')+e+e'$. A fixed dyadic schedule therefore yields a regular
name. Its upper-bound and approximation properties follow by comparison
with finer nets. No exact maximizer is selected.

For expectation, first replace the continuous integral by its quadrature,
costing $Lh+B\tau$ by Lemma~\ref{lem:tent-quadrature}. Replacing the
exact values at the nodes by the finite table costs at most $e$ because
the rational weights are nonnegative and normalized. This proves the
integration rule and its parameter modulus.

There are finitely many syntax nodes. The root bound is a finite
nonnegative expression in the local errors and mesh sizes. Positive
rational overestimates of its coefficients let us choose dyadic errors
small enough by bounded binary-length calculations. Each prescribed net
and mass sampler has its stated PR clock; composing them yields the
overall PR routine. For instance, make every local contribution at most
$2^{-p}/(1+s)$ after accounting for all its root influence coefficients,
where $s$ is the number of nodes. The computed bound then tends to zero
with an explicit rate. The same estimates comparing two runs give the
regularity of the resulting value name and complete the induction.
\end{proof}

If at an expectation node the table error varies with the cell, the more
accurate rule is
$\sum_iw_i e_i+Lh+B\tau$. Its soundness is
Lemma~\ref{lem:weighted}. Shared formula tables can be cached, with
their incoming influences accumulated by
Proposition~\ref{prop:adjoint}. This is where probability information
improves the numerical workload beyond a uniform worst-case schedule.

\subsection{Mixed norms and the order of binders}

The uniform rule is sufficient for every formula in the theorem. A
compiler can safely accept weaker hypotheses in special contexts.
If $\varphi(x,y)$ is followed by $\mathsf E_x$ and then $\sup_y$, an
appropriate error certificate is
\[
 \sup_y\int |\varphi(x,y)-\widetilde\varphi(x,y)|\,d\mu(x)\le e.
\]
If the order is $\sup_y$ first and $\mathsf E_x$ second, a sufficient
certificate is
\[
 \int\sup_y|\varphi(x,y)-\widetilde\varphi(x,y)|\,d\mu(x)\le e.
\]
The second is generally stronger. Their sufficiency follows from the
elementary inequalities for integration and supremum; they cannot be
interchanged merely because both contain an integral and a supremum.
Similarly, small product-space $L^1$ error by itself does not control
either supremum without a further uniformity certificate.

\begin{example}[Two genuinely different formulas]\label{ex:binder-order}
For $x,y\in[0,1]$ with Lebesgue probability measure,
\[
 \sup_y\E_x|x-y|=\frac12,
 \qquad
 \E_x\sup_y|x-y|=\frac34.
\]
Indeed $\E_x|x-y|=y^2/2+(1-y)^2/2$, whose maximum on $[0,1]$ is
$1/2$. On the other hand $\sup_y|x-y|=\max(x,1-x)$, with integral
$3/4$. On the $m$ equally weighted midpoint nodes, for even $m$, the
two finite values are respectively
\[
 \frac12-\frac1{2m},\qquad\frac34-\frac1{2m}.
\]
For the first, the extreme midpoint maximizes a convex average of
absolute values, and direct summation gives $(m-1)/(2m)$. For the second,
the maximum is $\max(x,1-x)-1/(2m)$ on each midpoint; the symmetric
midpoint sum of $|x-1/2|$ is $1/4$. Thus the finite evaluator preserves
the nonzero gap exactly. The supplied reference program computes both
syntax trees independently.
\end{example}

Two expectation binders over independent product measures can be
exchanged, with proof by finite Fubini followed by the named limit.
Their exchange is a proved rule, not a general permission to reorder
binders. The example shows why the syntax must retain the difference.

\subsection{Model-theoretic boundaries}

The metric space of all events of an atomless probability algebra,
with distance $\mu(A\mathbin\triangle B)$, need not be totally bounded.
Independent fair-coordinate events have pairwise distance $1/2$.
Thus the compact-net theorem does not evaluate unrestricted quantifiers
over all events. It evaluates compact point sorts, supplied compact
families of observables, or explicitly finite event presentations.
The model theory of probability algebras has substantially richer
structure \cite{berenstein-henson}; its full quantifier elimination and
model existence statements are not consequences of this finite theorem.

A computed value interval can certify a strict inequality separated by
a rational margin. It cannot in general decide exact zero truth.
Similarly, values of every finite test do not by themselves select one
coherent model of an arbitrary theory. For that task one needs compatible
named structure data or an additional constructive extension procedure.
These are mathematical boundaries of the embedding, not defects that can
be removed by replacing words with certificates.
